\subsection{研究对象与预测任务}\label{sec:forecast-task}
经验研究对象是历史GPX集合中带保存时间的地理位置记录。
每条记录描述一个身份未知的记录携带者的位置变化；保留档案不能确认其是在步行、骑行
还是使用其他交通方式，也不能确认两条记录是否属于同一个人。
因此本文研究入选记录中的位置运动，不把它解释为已定义参与者人群的行为规律。
清理段是从记录派生的一段序列；窗口将该段分为已观测前缀与未来目标。
同一个记录哈希块可以包含多个窗口，主评估只从每个入选块选取一个预测起点。

对保存时间为$t_0$的起点，令$\mathcal H_0$包含允许可见的位置／时间前缀，
末点为当前位置$x_0$，最后三对用于初始化保留运动历史缓冲。
任务是给定这些信息及开发数据拟合后固定的参数$\widehat\xi$，预测未来位置。
方法分支内部东西／南北坐标系以整个前缀的首个可见位置为锚，未必是最后三点中的首点；
地形预测和评价则以末个可见位置，即预测起点为锚。
方法输出经地理坐标转换到统一评分坐标系，式\eqref{eq:forecast-object}的分布以该坐标表示。
方法分支还读取保留时钟及其太阳代理，地形分支另外在模拟位置查询固定地图场。
真实未来位置用于队列筛选和评价，不用于更新预测。
对四个实际匹配的经过时长$u_1,\ldots,u_4$，报告对象为
\begin{equation}\label{eq:forecast-object}
 \widehat P_j=\mathcal L_{\widehat\xi}(X_{t_0+u_j}\mid\mathcal H_0),
 \qquad
 \widehat P_{N,j}=\frac1N\sum_{p=1}^{N}\delta_{X_j^{(p)}},
 \quad j=1,\ldots,4 .
\end{equation}
其中$\delta_x$是位置$x$处的点质量，同一模拟粒子$p$在四个目标时刻各有一个位置。
粒子集合近似模型规定的分布，粒子数不是人数；集合均值给出代表性位置，
径向分位数给出预测圆。这两种摘要与整个位置分布分别回答不同问题。
本研究没有另行拟合观测误差模型：保存目标位置作为评价参照，其物理精度仍未独立核实。

开发窗口用于估计参数，最终评估窗口提供预测起点和真实结果。
主分析考察固定预测器对这些结果的描述能力怎样随保存预测时域变化；
它不估计人类运动的因果定律，也不估计实际搜索中找到某人的概率。
